For real phase retrieval, the critical number of measurements is \(N=2m-1\)
\cite{BalanCasazzaEdidin2006,BandeiraCahillMixonNelson2014}.  When \(A\in\R^{(2m-1)\times m}\) is full
spark, the complement property
\cite{BalanCasazzaEdidin2006,BandeiraCahillMixonNelson2014,ConcaEdidinHeringVinzant2015} implies that
the Balan--Wang parameter \cite{BalanWang2015} reduces to the extremal
square-row-submatrix quantity:
\[
  \omega(A)=\min_{|T|=m}\smin(A_T).
\]

\begin{theorem}[Continuous critical frames: lower side]\label{thm:pr-lower}
Let \(A_m\in\R^{(2m-1)\times m}\) have independent entries with densities
uniformly bounded by \(K\).  Then for every \(\varepsilon>0\),
\[
  \PP\{\omega(A_m)\le e^{-(\log4+\varepsilon)m}\}\to0,
\]
provided \(A_m\) is interpreted on the full-spark event.  For continuous entry
laws the full-spark event has probability one.
\end{theorem}

\begin{proof}
At \(N=2m-1\), Stirling's formula gives
\[
  \frac{1}{m}\log \binom{2m-1}{m}\to\log4.
\]
Apply Corollary~\ref{cor:bd-lower}.  Continuous iid entries give full spark
almost surely; more generally, independent absolutely continuous entries have
this property because each \(m\times m\) minor is a nonzero polynomial in the
entries, and there are only finitely many minors for each \(m\).
\end{proof}

\begin{corollary}[Directional anti-concentration critical frames]\label{cor:pr-directional-lower}
Let \(A_m\in\R^{(2m-1)\times m}\) have independent rows satisfying the affine
directional small-ball hypothesis of Corollary~\ref{cor:directional-row-lower}
with constants \(L_{m,i}\) such that
\[
  \frac1m\log\left(\sum_{i=1}^{2m-1}L_{m,i}\right)\to0.
\]
Then \(A_m\) is full spark almost surely and, for every \(\varepsilon>0\),
\[
  \PP\{\omega(A_m)\le e^{-(\log4+\varepsilon)m}\}\to0
\]
under the Balan--Wang singular-value convention used here.
\end{corollary}

\begin{proof}
For each fixed \(m\), the directional small-ball hypothesis has no atoms on
affine hyperplanes.  Conditioning row by row shows that every selected
\(m\times m\) minor is nonsingular almost surely, and a finite union over
minors gives full spark.  At \(N=2m-1\), full spark identifies
\(\omega(A_m)\) with \(\min_{|T|=m}\smin(A_{m,T})\).  Apply
Corollary~\ref{cor:directional-row-lower} with \(h(2)=\log4\).
\end{proof}

\begin{corollary}[Conditional universal Balan--Wang base]\label{cor:pr-cond}
If, in addition, the real bounded-density ensemble satisfies the
pair-correlation condition of Section~\ref{sec:universality} at \(\gamma=2\),
then
\[
  \frac{1}{m}\log\omega(A_m)\xrightarrow{\PP}-\log4,
  \qquad
  \omega(A_m)=4^{-m+o_\PP(m)}.
\]
\end{corollary}

The Gaussian source paper \cite{ShmaloExtremal2026} proves the hypothesis in
the Gaussian case and hence answers the exponential-base part of Randomstrasse
Open Problem 21~\cite{BandeiraStablePhase} for iid Gaussian measurements.
Appendix~\ref{sec:app-keystone} establishes the pair-correlation condition for
all iid bounded-density real ensembles by an elementary off-support resolvent
estimate, self-contained (Remark~\ref{appB:selfcontained}) with the anisotropic
local law optional, so the universal base \(\omega(A_m)=4^{-m+o_\PP(m)}\) holds
unconditionally for those ensembles; sharpening the rate and extending beyond
bounded density is Problem~\ref{prob:pc}.

\begin{corollary}[Row-normalized Gaussian critical frames]\label{cor:pr-normalized-gaussian}
Let \(G_m\in\R^{(2m-1)\times m}\) have iid standard Gaussian entries and let
\(U_m\) be obtained by normalizing each row of \(G_m\) to Euclidean norm one.
Then \(U_m\) is full spark almost surely and its Balan--Wang singular-value
statistic satisfies
\[
  \frac1m\log\omega(U_m)\xrightarrow{\PP}-\log4,
  \qquad
  \omega(U_m)=4^{-m+o_\PP(m)}.
\]
\end{corollary}

\begin{proof}
Row normalization rescales each nonzero Gaussian row and therefore does not
change which row submatrices are singular.  Since \(G_m\) is full spark almost
surely and has no zero row almost surely, \(U_m\) is full spark almost surely.
At \(N=2m-1\), full spark identifies the Balan--Wang statistic with
\(\min_{|T|=m}\smin((U_m)_T)\).  Corollary~\ref{cor:row-normalized-gaussian}
with \(\gamma=2\) gives the limit because \(h(2)=\log4\).
\end{proof}

The same conclusion holds for iid rows uniform on any sphere of radius
\(e^{o(m)}\), by Corollary~\ref{cor:spherical-pr}; the change from the unit
sphere to radius \(\sqrt m\), for example, is only a scalar factor
\(\sqrt m=e^{o(m)}\).

\begin{corollary}[Diagonal preconditioning of Gaussian critical frames]\label{cor:pr-diagonal-preconditioned}
Let \(G_m\in\R^{(2m-1)\times m}\) have iid standard Gaussian entries.  Let
\(E_m=\operatorname{diag}(e_{1m},\ldots,e_{mm})\) have strictly positive finite
random diagonal entries, possibly dependent on \(G_m\), and assume
\[
  \max_{1\le j\le m}|\log e_{jm}|=o_\PP(m).
\]
Put \(B_m=G_mE_m\).  Then \(B_m\) is full spark almost surely and
\[
  \frac1m\log\omega(B_m)\xrightarrow{\PP}-\log4,
  \qquad
  \omega(B_m)=4^{-m+o_\PP(m)}.
\]
In particular, the same base holds when the columns of \(G_m\) are normalized
to unit Euclidean norm.
\end{corollary}

\begin{proof}
Right multiplication by an invertible diagonal matrix multiplies every
\(m\times m\) row minor determinant by \(\det E_m\), so full spark is preserved
from \(G_m\) to \(B_m\).  On the full-spark event,
\[
  \omega(B_m)=\min_{|T|=m}\smin((G_m)_T E_m).
\]
With \(\Delta_m=\max_j|\log e_{jm}|\), Proposition~\ref{prop:random-column-weight}
gives
\[
  e^{-\Delta_m}\omega(G_m)\le \omega(B_m)\le e^{\Delta_m}\omega(G_m).
\]
The Gaussian Balan--Wang law for \(\omega(G_m)\), together with
\(\Delta_m=o_\PP(m)\), proves the claim.  The column-normalized case is
Corollary~\ref{cor:column-normalized-gaussian} at \(\gamma=2\).
\end{proof}

\begin{remark}
Column diagonal scaling is a coordinate preconditioning of the unknown signal,
not normalization of the measurement vectors.  It preserves real
phase-retrieval injectivity because the preconditioner is invertible, and it
preserves the exponential Balan--Wang base only when the diagonal entries
themselves have subexponential logarithmic extrema.  A subexponential condition
number alone is not enough, since a scalar factor \(e^m I\) shifts the absolute
singular-value exponent.
\end{remark}

\begin{remark}[Supercritical measurement counts]
The random-row-count theorem in Section~\ref{sec:universality} is an RMT
statement about the square statistic \(\min_{|T|=m}\smin(A_T)\).  It is not a
Balan--Wang corollary when the measurement count merely fluctuates around
\(2m\).  For a fixed supercritical count \(N=2m-1+r\), the fixed-excess quantity
\[
  \min_{|T|=m+r}\smin(A_T),
\]
is the natural Balan--Wang \(\omega\)-branch under the singular-value convention
of \cite{BalanWang2015}.  It should not be confused with the full lower
Lipschitz partition statistic, where partitions whose two sides both span can
also enter.  The fixed-excess results of Section~\ref{sec:geometric} should
therefore be read as RMT input for supercritical stability, not as a theorem
about every Balan--Wang stability convention.
\end{remark}

\begin{corollary}[Directional fixed-excess branch]\label{cor:pr-directional-branch}
Fix \(r\ge1\), and let \(A_m\in\R^{(2m-1+r)\times m}\) have independent rows
which satisfy the affine directional small-ball hypothesis of
Corollary~\ref{rect:directional-row-fixed-excess} with constants \(L_{m,i}\).
Define the fixed-excess branch statistic
\[
  \omega_r^{\rm fe}(A_m)
  =
  \min_{\substack{T\subset[2m-1+r]\\ |T|=m+r}}\smin((A_m)_T).
\]
Assume
\[
  \frac1m\log
  \sum_{\substack{S\subset[2m-1+r]\\ |S|=r+1}}
  \prod_{i\in S}L_{m,i}\longrightarrow0 .
\]
Then, for every \(\varepsilon>0\),
\[
  \PP\left\{\omega_r^{\rm fe}(A_m)
  \le \exp\left[-\left(\frac{\log4}{r+1}+\varepsilon\right)m\right]\right\}
  \to0 .
\]
On any common coupling the same hypotheses give the almost-sure lower envelope
\[
  \liminf_{m\to\infty}\frac1m\log\omega_r^{\rm fe}(A_m)
  \ge -\frac{\log4}{r+1}
  \qquad\hbox{almost surely}.
\]
This is the deficient-complement branch of the Balan--Wang partition geometry.
For \(r>0\) it is not the full lower-Lipschitz partition statistic.
\end{corollary}

\begin{proof}
Set \(p=r+1\) and \(N_m=2m-1+r\).  In
Corollary~\ref{rect:directional-row-fixed-excess},
\[
  \binom{N_m-p}{m-1}=\binom{2m-2}{m-1},
  \qquad
  \frac1m\log\binom{2m-2}{m-1}\to\log4 .
\]
The assumed elementary-symmetric directional budget is subexponential, and the
polynomial slack \((p/2)\log m\) is \(o(m)\).  Taking
\(u_m=p\varepsilon m\) gives the convergence-in-probability estimate.

For the almost-sure lower envelope, take \(u_m=2\log m\) in the same finite
budget estimate and apply Borel--Cantelli at every fixed \(\varepsilon>0\).
The directional small-ball hypothesis also implies full spark almost surely in
each fixed dimension by conditioning row by row, but full spark does not turn
the displayed branch into the full supercritical stability statistic.
\end{proof}

\begin{corollary}[Almost-sure Balan--Wang lower envelope]\label{cor:pr-as-lower-envelope}
Let \((A_m)_{m\ge1}\) be real \((2m-1)\times m\) random matrices on a common
probability space, with independent entries inside each \(A_m\) and Lebesgue
densities bounded by a common constant.  On the simultaneous full-spark event,
\[
  \liminf_{m\to\infty}\frac1m\log\omega(A_m)\ge-\log4
  \qquad\hbox{almost surely}.
\]
For continuous entry laws the simultaneous full-spark event has probability
one.  The same almost-sure lower envelope holds for the row-normalized
Gaussian critical frames \(U_m\) of
Corollary~\ref{cor:pr-normalized-gaussian}.  It also holds for
diagonally preconditioned Gaussian critical frames \(B_m=G_mE_m\) whenever
\(\max_j|\log e_{jm}|=o(m)\) almost surely; in particular it holds for
column-normalized Gaussian critical frames.
\end{corollary}

\begin{proof}
The first statement is Corollary~\ref{cor:as-square-lower-envelope} at
\(\gamma=2\), where \(h(2)=\log4\), combined with the full-spark reduction
\(\omega(A_m)=\min_{|T|=m}\smin(A_{m,T})\).  Simultaneous full spark holds
almost surely because the exceptional event is a countable union, over all
\(m\) and all selected minors, of zero sets of nonzero determinant polynomials.

For normalized Gaussian frames, row normalization preserves full spark and
\[
  \omega(U_m)\ge
  \frac{\omega(G_m)}{\max_{i\le 2m-1}\|(G_m)_i\|_2}.
\]
The Gaussian norm tail gives
\[
  \sum_m
  \PP\left\{\max_{i\le2m-1}\|(G_m)_i\|_2>e^{\delta m}\right\}<\infty
  \qquad(\delta>0).
\]
Combining this with the Gaussian lower envelope for \(\omega(G_m)\) and then
letting \(\delta\downarrow0\) proves the normalized statement.

For \(B_m=G_mE_m\), the proof of
Corollary~\ref{cor:pr-diagonal-preconditioned} gives
\[
  \omega(B_m)\ge e^{-\Delta_m}\omega(G_m),
  \qquad \Delta_m=\max_j|\log e_{jm}|.
\]
If \(\Delta_m=o(m)\) almost surely, the Gaussian lower envelope transfers.
The column-normalized case has this property by the summable column-norm
estimates in Corollary~\ref{cor:column-normalized-gaussian}.
\end{proof}

\begin{remark}
The corollary concerns the Balan--Wang singular-value statistic \(\omega\).
Other stability conventions may square singular values or include additional
nonexponential factors; those conventions have the corresponding transformed
exponential base.  Corollary~\ref{cor:pr-as-lower-envelope} is only a lower
envelope; it does not upgrade the exact Gaussian base to almost-sure
convergence.
\end{remark}

\begin{remark}[Discrete frames]
For atomic entries, full spark is no longer automatic.  Proposition
\ref{prop:finite-alphabet-obstruction} gives deterministic finite-alphabet
projective-pair full-spark obstructions at the corresponding row count.  In the
real sign case, Proposition~\ref{prop:sign-critical} gives the critical count
exactly: every real sign matrix with \(2m-1\) rows and \(m\ge2\) columns has a
singular \(m\)-row submatrix.  Thus iid Rademacher critical frames fail
injectivity with probability one by the complement property
\cite{BalanCasazzaEdidin2006,BandeiraCahillMixonNelson2014}, in sharp contrast
with continuous Gaussian frames.  Away from the critical real sign case,
including complex finite alphabets, a singular \(m\)-row submatrix is only a
full-spark or square-minor obstruction and is not claimed to settle
phase-retrieval injectivity.
\end{remark}

\begin{problem}[Complex phase-retrieval stability base]
The complex injectivity threshold and the complex stability geometry are not
obtained by simply replacing \(d=1\) with \(d=2\) at \(N=2m-1\).  At the generic
complex threshold \(4m-4\) \cite{ConcaEdidinHeringVinzant2015}, determine the
correct Balan--Wang-type exponential base for complex random measurements, and
relate it to extremal singular values of the appropriate realified submatrix
families.
\end{problem}
